\subsection{An explicit six-regime allocation diagram}\label{sec:strata-phase}
Here we solve the variational comparator in Theorem~\ref{thm:stratified}, rather than claim that its implicit comparison constants are one. Take two components with $(d_1,\alpha_1,\lambda_1)=(1,1,1/2)$ and $(d_2,\alpha_2,\lambda_2)=(2,2,1/2)$, and total dictionary width $m=4$. If $s$ is the selected mass from the first component, the comparator is
\[
 J_b(s)=\max\left\{\frac{4s^3}{b^2},\frac{s^2}{2b}\right\}
       +\max\left\{\frac{2(1-s)^2}{b},
                     \frac{(1-s)^{3/2}}{\sqrt{2b}}\right\},
\]
with $\max(0,1-b/2)\le s\le\min(1,b/2)$ and $b\ge1$. Its minimum describes the risk up to the inherited component-dependent constants. It is not an exact formula for every acquired law having these exponents.

These component hypotheses can arise from an actual bounded preparation: take the first location uniform on a unit segment in $\R^2$ and the second uniform on a unit square, with reliability variables having distribution functions $t$ and $t^2$, respectively. With a fresh fair sign $E$, audit $Z=Y+TEe_1$ and report $(Y,T)$. Then $f=Y$ and $w=T^2$ in a common compact convex carrier. The mixture has the stated acquired dimensions and exponents. The diagram still concerns the theorem's comparator rather than replacing its comparison constants by exact ones.

\begin{proposition}[Exact selected-mass phases for the comparator]\label{prop:strata-phases}
Put
\[
 b_1=6-2\sqrt6,\quad b_2=(16\sqrt6-24)/5,\quad
 b_3=8(\sqrt2-1),\quad b_4=32/5,\quad b_5=36/5.
\]
The minimizer is unique and continuous in $b$. It is
\begin{equation}\label{eq:strata-phases}
 s_*(b)=
 \begin{cases}
 1-b/2,&1\le b\le b_1,\\
 (\sqrt{b^2+12b}-b)/6,&b_1\le b\le b_2,\\
 1-b/8,&b_2\le b\le b_3,\\
 \text{the unique }s\in(0,1)\text{ with }128s^4=b^3(1-s),
       &b_3\le b\le b_4,\\
 b/8,&b_4\le b\le b_5,\\
 (\sqrt{81b^2+288b}-9b)/16,&b\ge b_5.
 \end{cases}
\end{equation}
In particular there are two genuine intervals where the minimizer lies at a nonsmooth crossover, not just isolated changes between differentiable stationary formulas. The initially decreasing allocation to the first component eventually increases to one.
\end{proposition}
\begin{proof}
Each summand is strictly convex on its nonnegative mass interval, and hence $J_b$ is strictly convex. Both components change from their location-limited term to their reliability-limited term when their own selected mass crosses $b/8$. On a differentiable interval the minimizing equation is equality of their marginal costs.

When both are reliability-limited, this equation is
$12s^2/b^2=4(1-s)/b$, or $3s^2=b(1-s)$. Its solution first enters the feasible interval at $s=1-b/2$, giving $b_1$, and reaches the second component's crossover $1-s=b/8$ at $b_2$. At that crossover the second marginal-cost interval is $[3/8,1/2]$. The first marginal cost $12(1-b/8)^2/b^2$ decreases through that interval as $b$ runs from $b_2$ to $b_3$. Thus $s=1-b/8$ is the minimizer throughout that interval.

With the first component reliability-limited and the second location-limited, marginal equality is
$12s^2/b^2=3\sqrt{1-s}/(2\sqrt{2b})$, equivalent to the displayed quartic. Its left side is strictly increasing in $s$ and its right side strictly decreasing. The root reaches $s=b/8$ at $b_4$. At this first-component crossover its own marginal-cost interval is $[1/8,3/16]$. The second marginal cost $3\sqrt{1-b/8}/(2\sqrt{2b})$ decreases from $3/16$ to $1/8$ between $b_4$ and $b_5$, giving the second nonsmooth interval. Thereafter both terms are location-limited; equality becomes $8s^2=9b(1-s)$, with the final displayed root. Substitution verifies the active conditions and the original mass caps in each interval. Strict convexity proves uniqueness; matching endpoint values proves continuity. The last equation implies $1-s_*(b)\sim8/(9b)$.
\end{proof}

\begin{figure}[htbp]
\centering
\setlength{\unitlength}{1pt}
\begin{picture}(410,210)(0,-14)
\put(38,22){\line(1,0){350}}
\put(38,22){\line(0,1){145}}
\put(0,182){$s_*(b)$}
\put(34,22.000){\line(1,0){4}}
\put(8,19.000){\makebox(23,8)[r]{$ 0 $}}
\put(34,93.000){\line(1,0){4}}
\put(8,90.000){\makebox(23,8)[r]{$ 0.5 $}}
\put(34,164.000){\line(1,0){4}}
\put(8,161.000){\makebox(23,8)[r]{$ 1 $}}
\put(38.000,18){\line(0,1){4}}
\put(30.000,3){\makebox(16,9){$ 1 $}}
\put(109.328,18){\line(0,1){4}}
\put(101.328,3){\makebox(16,9){$ 2 $}}
\put(203.619,18){\line(0,1){4}}
\put(195.619,3){\makebox(16,9){$ 5 $}}
\put(274.947,18){\line(0,1){4}}
\put(266.947,3){\makebox(16,9){$ 10 $}}
\put(388.000,18){\line(0,1){4}}
\put(380.000,3){\makebox(16,9){$ 30 $}}
\put(174,-11){$b$ (logarithmic scale)}
\thinlines
\put(47.903,22){\line(0,1){142}}
\put(40.903,170.000){\makebox(14,8){$b_1$}}
\thinlines
\put(152.360,22){\line(0,1){142}}
\put(145.360,178.000){\makebox(14,8){$b_2$}}
\thinlines
\put(161.287,22){\line(0,1){142}}
\put(154.287,170.000){\makebox(14,8){$b_3$}}
\thinlines
\put(229.022,22){\line(0,1){142}}
\put(222.022,178.000){\makebox(14,8){$b_4$}}
\thinlines
\put(241.143,22){\line(0,1){142}}
\put(234.143,170.000){\makebox(14,8){$b_5$}}
\thicklines
\qbezier(38.000,93.000)(40.734,91.062)(43.469,89.125)
\qbezier(43.469,89.125)(45.686,87.476)(47.903,85.828)
\qbezier(47.903,85.828)(48.420,85.942)(48.938,86.055)
\qbezier(48.938,86.055)(51.672,86.661)(54.406,87.267)
\qbezier(54.406,87.267)(57.141,87.878)(59.875,88.489)
\qbezier(59.875,88.489)(62.609,89.105)(65.344,89.720)
\qbezier(65.344,89.720)(68.078,90.341)(70.812,90.961)
\qbezier(70.812,90.961)(73.547,91.585)(76.281,92.210)
\qbezier(76.281,92.210)(79.016,92.838)(81.750,93.466)
\qbezier(81.750,93.466)(84.484,94.098)(87.219,94.730)
\qbezier(87.219,94.730)(89.953,95.365)(92.688,96.001)
\qbezier(92.688,96.001)(95.422,96.639)(98.156,97.277)
\qbezier(98.156,97.277)(100.891,97.918)(103.625,98.558)
\qbezier(103.625,98.558)(106.359,99.201)(109.094,99.844)
\qbezier(109.094,99.844)(111.828,100.488)(114.562,101.133)
\qbezier(114.562,101.133)(117.297,101.779)(120.031,102.424)
\qbezier(120.031,102.424)(122.766,103.071)(125.500,103.718)
\qbezier(125.500,103.718)(128.234,104.365)(130.969,105.013)
\qbezier(130.969,105.013)(133.703,105.660)(136.438,106.307)
\qbezier(136.438,106.307)(139.172,106.954)(141.906,107.601)
\qbezier(141.906,107.601)(144.641,108.247)(147.375,108.893)
\qbezier(147.375,108.893)(149.868,109.481)(152.360,110.069)
\qbezier(152.360,110.069)(152.602,109.942)(152.844,109.815)
\qbezier(152.844,109.815)(155.578,108.336)(158.312,106.858)
\qbezier(158.312,106.858)(159.800,106.020)(161.287,105.182)
\qbezier(161.287,105.182)(162.534,105.742)(163.781,106.301)
\qbezier(163.781,106.301)(166.516,107.538)(169.250,108.774)
\qbezier(169.250,108.774)(171.984,110.020)(174.719,111.266)
\qbezier(174.719,111.266)(177.453,112.519)(180.187,113.772)
\qbezier(180.187,113.772)(182.922,115.029)(185.656,116.286)
\qbezier(185.656,116.286)(188.391,117.543)(191.125,118.800)
\qbezier(191.125,118.800)(193.859,120.055)(196.594,121.309)
\qbezier(196.594,121.309)(199.328,122.556)(202.062,123.803)
\qbezier(202.062,123.803)(204.797,125.039)(207.531,126.275)
\qbezier(207.531,126.275)(210.266,127.496)(213.000,128.716)
\qbezier(213.000,128.716)(215.734,129.916)(218.469,131.116)
\qbezier(218.469,131.116)(221.203,132.292)(223.938,133.467)
\qbezier(223.938,133.467)(226.480,134.534)(229.022,135.600)
\qbezier(229.022,135.600)(229.214,135.812)(229.406,136.025)
\qbezier(229.406,136.025)(232.141,139.137)(234.875,142.248)
\qbezier(234.875,142.248)(237.609,145.530)(240.344,148.812)
\qbezier(240.344,148.812)(240.743,149.306)(241.143,149.800)
\qbezier(241.143,149.800)(243.478,150.060)(245.812,150.319)
\qbezier(245.812,150.319)(248.547,150.614)(251.281,150.908)
\qbezier(251.281,150.908)(254.016,151.192)(256.750,151.476)
\qbezier(256.750,151.476)(259.484,151.750)(262.219,152.024)
\qbezier(262.219,152.024)(264.953,152.287)(267.688,152.551)
\qbezier(267.688,152.551)(270.422,152.805)(273.156,153.059)
\qbezier(273.156,153.059)(275.891,153.303)(278.625,153.548)
\qbezier(278.625,153.548)(281.359,153.783)(284.094,154.018)
\qbezier(284.094,154.018)(286.828,154.244)(289.562,154.470)
\qbezier(289.562,154.470)(292.297,154.687)(295.031,154.904)
\qbezier(295.031,154.904)(297.766,155.112)(300.500,155.320)
\qbezier(300.500,155.320)(303.234,155.520)(305.969,155.720)
\qbezier(305.969,155.720)(308.703,155.912)(311.437,156.103)
\qbezier(311.437,156.103)(314.172,156.287)(316.906,156.471)
\qbezier(316.906,156.471)(319.641,156.647)(322.375,156.823)
\qbezier(322.375,156.823)(325.109,156.992)(327.844,157.161)
\qbezier(327.844,157.161)(330.578,157.322)(333.312,157.484)
\qbezier(333.312,157.484)(336.047,157.638)(338.781,157.793)
\qbezier(338.781,157.793)(341.516,157.940)(344.250,158.088)
\qbezier(344.250,158.088)(346.984,158.230)(349.719,158.371)
\qbezier(349.719,158.371)(352.453,158.506)(355.188,158.641)
\qbezier(355.188,158.641)(357.922,158.770)(360.656,158.899)
\qbezier(360.656,158.899)(363.391,159.022)(366.125,159.146)
\qbezier(366.125,159.146)(368.859,159.263)(371.594,159.381)
\qbezier(371.594,159.381)(374.328,159.494)(377.062,159.606)
\qbezier(377.062,159.606)(379.797,159.713)(382.531,159.820)
\qbezier(382.531,159.820)(385.266,159.923)(388.000,160.025)
\end{picture}

\caption{The exact minimizer of the displayed comparator, not a fitted risk curve. Vertical lines mark the five algebraic transition budgets. The two crossover intervals are part of the optimum; smoothing a maximum before differentiation would miss them. The plotted values are checked against independent bounded scalar minimization.}\label{fig:strata-phase}
\end{figure}
